% !TEX root = ../main.tex
\chapter{Architectural Design}

The architecture follows mainly from two requirements: the non-functional requirements that asks for model updates not
to depend on anything other than the state they are given and the time elapsed, in order to allow the game behavior
to be exercised without a window ever being opened, and the implementation requirement that asks that the rules do
not depend on the chosen graphical library. The obtained architecture is one in which everything that determines what
happens in a game, such as how far the player and the enemies advance in the time that has passed, which cell an enemy
goes towards next, what a collected item is worth, is computed from the state of the game together with a duration,
and yields the state that follows. These actions are pure: none of them opens a file, reads the system clock, or
touches the screen. The one place in which the model does consult the clock is when a finished game is turned into a
result, which is stamped with the moment it was achieved, but that is not part of any update.
The model does not contain any mutable variable at all: the few mutable variables are all in the entry point and the
interface.

The three things that cannot be computed that way are kept out of that computation altogether: the window and the
drawing done upon it, the timer that reports when a frame is due together with the keys the player has passed, and
the reading and writing of mazes, saved games and recorded results. They are confined to a layer that is around the
computation: the entry point, that owns the window and the timer, and adapters that own the files.

\section{Architectural patterns}

\subsection{Functional core, imperative shell}

The system is divided into a \textbf{core} that is entirely pure and a \textbf{shell} that is not. The core holds
every rule of the game: given the state of a game and the time that has passed, it computes the next state and
returns it. It performs no input or output, reads no clock, and holds no mutable variable. The shell is the part
that must be impire: the graphical library, the frame timer, the file system.

Nearly all mutability in the running application is confined to a single reference at the boundary, which the frame
timer replaces with the state the core has computed. Everything below it is immutable, and that allows the game to be
driven from a test at any speed, replayed or even suspended and resumed.

\subsection{Model View Controller}

Within the division into core and shell, the MVC architecture pattern is applied. The \textbf{model} holds the domain
and its rules, while not concerning itself with the rest. The \textbf{view} draws and collects input, and contains no
rule of its own. The \textbf{controller} sits between them, translating what the player does into changes of state and
what the model becomes into something the view can draw. Each update produces a \textbf{level view}, which is a flat
description of what should appear on screen, that is positions, lives, score, remaining time, effects, and it is what
the view receives. SO the view cannot reach into the game's state, decoupling it completely from the model.

\subsection{Ports and adapters at the boundary}

Everything that is outside the program, such as the mazes it offers, the mazes a player has written, saved games,
recorded results, and the remembered player name, is reached only through interfaces that describe what is wanted
rather than how it is stored, satisfying the implementation requirement that asks that the model be independent of
the interface and the graphical library.

The knowledge of storage is concentrated: on class knows the layout of the folders, one knows how a saved game is
written, one knows how a result is encoded. Each of them can be replaced without anything else in the program being
aware, and each is replaced during testing, where what would normally be outside the game is supplied by a substitute
that records what was asked of it and can be told to fail on demand. This is what allows the controller to be tested
with no file ever being written.

The separation is not total. Three of these operations take a file path, because everything the game remembers is kept
in plain files under a folder of the player, rather than in a database, and because the player picks a maze or a saved
game through a file dialog, which yields a path.

\section{Main components}

\begin{description}
  \item[Entry point] The \textbf{window}, together with the two sources of activity that drive the program: the keys
  that are pressed, and a timer that fires once per frame. It builds the window, installs the\textbf{ key filter}
  and starts the \textbf{timer}, and holds the reference to the current state of the application, replacing it
  with whatever each of those two events produces.  On a key press, it hands the key to the input mapping and
  the direction or command that comes back to the application. On a frame, it hands the application that frame's
  timestamp, and receives the application as it then stands, which it draws. When a game ends, it shows the menu
  on screen.
  \item[Application] is the state of the program as a whole: whether a menu is showing or a game is being played,
  and what the player was last told. It interprets \textbf{commands} that the interface raises, asks for the
  necessary resources, and decides when a game result is to be recorded in the leaderboard. It exchanges commands and
  application states with the entry point, and levels with the game session.
  \item[Game session] is a game in progress, holding everything about it that changes from one frame to the next:
  the level as it now stands, the rendering through which the interface is notified of it, whether the loop
  is running or paused, and the moment of the last frame it was advanced to.  Given the moment of a new frame
  it yields the session as it stands at that moment, and the interval between the two becomes the duration it
  hands to the model's update pipeline, never negative and never more than fifty milliseconds, so that a frame
  the machine was slow over cannot carry everyone across the maze at once with every collision along the way.
  Each advance notifies the interface on the level that resulted. A session that is paused is left untouched.
  \item[Environment]. Everything outside the program is reached through interfaces owned by the code that needs them
  rather than by the code that implements them. Two are ports in the strict sens, each naming a single capability
  with no mention of how it is realised.
  \item[View] is the graphical interface: the \textbf{menu} for choosing a name, a maze and a mode, the \textbf{board}
  on which a game is drawn, the \textbf{overlay} shown when a game is paused or has ended, the \textbf{window}
  listing recorded results. It receives\textbf{ level views} and raises, from its buttons, the choices a player
  makes, that are starting, loading, pausing, restarting, saving and returning to the menu.
  \item[Level and update pipeline]. This is the heart of the core. A level is the whole state of a game at one instant.
  The pipeline is the fixed sequence of stages, each a function from one state to the next, run in a fixed order:
  input, enemy decisions, movement, collisions and teleports, gathering, effects, and a final reconciliation. It
  is generic in the state it carries, so the order is fixed once and says nothing about what level it is. A level
  supplies the stages that give each step its meaning, with the duration to advance by already built into them,
  leaving a tick a plain function from a level to the level that follows.
  \item[Entities and movement] The player and the enemies, their positions, the directions they face and the crossings
  they are part-way through. They are used by the pipeline.
  \item[Enemy behavior] The strategies by which each kind of enemy chooses where to go, and the stage that applies them
  to every idle enemy. It receives a level and returns the enemies with their next movements decided.
  \item[Input mapping] is the single place where a key stands for an action: the arrow keys and W, A, S and D for
  the four directions, and escape for pausing and resuming. It is consulted by the entry point, which owns the
  key filter and passes on whatever comes back. Keeping it apart is what lets the controls be changed in one place.
  \item[Rendering] turns a level into the flat description the interface draws, and tells whoever is listening when
  that description changes. It holds the projection, the listeners registered against it, and the last
  description shown. Given a level it computes the new description and passes it on only if it differs from the
  last, so a frame in which nothing visible changed is not redrawn. It receives levels from the game session and
  hands \textbf{level views} to the view.
  \item[Collectibles and effects] is what a maze still has to gather and which temporary effects are in force. Gathering
  answers three things at once: what remains, the effect the item granted together with the instant it expires,
  and the points it was worth. The gathering and effect stages of the pipeline work upon it, collision resolution
  asks it whether invulnerability is is effect before deciding what a meeting costs, and rendering reads from it
  what is left and what is active.
  \item[Modes and scoring]  is what distinguishes the different ways of playing. The objective, the increasing
  difficulty, the points awarded for the various events, together with the score and the result that is left
  behind after a game is completed. After every update the level asks the mode, given the lives left, what remains
  to gather and the time elapsed, whether the game is still running, has been won or lost. The application asks
  for a result once the game is over, and hands it to the environment to record it.
  \item[Persistence adapters] The adapters behind the ports: saved games written as versioned property files carrying
  a copy of their maze, and leaderboards written as text through an encoding defined for a result.
  \item[Map pipeline] Reading a maze from a file, parsing its text into a grid, validating that grid against the rules
  that make a maze playable, and generating one where required. It receives text and returns either a validated
  maze or the reasons it was refused.
\end{description}
